\label{cap:introducao_modelo}

Este Trabalho de Conclusão de Curso (TCC) foi desenvolvido em parceria com Igor Lamoia Queiroz, aluno do Curso de Engenharia de Computação do CEFET-MG, Unidade Leopoldina. Os dois trabalhos tratam do mesmo projeto, mas com enfoques distintos, conforme a divisão de responsabilidades descrita no Apêndice~\ref{ap:plano}. Este texto descreve o núcleo do compilador: as fases de análise léxica, análise sintática, geração de código intermediário e execução pelo interpretador. O trabalho complementar, intitulado \textit{Interpretador Dinâmico para Ensino de Programação: Uma abordagem interativa com Ambiente de Desenvolvimento Integrado (IDE) para personalização e execução de comandos definidos pelo usuário}, foca na IDE web, na interface de personalização da linguagem e na experiência de uso da plataforma. Por se tratar de trabalhos complementares e interdependentes, para a compreensão do sistema como um todo, recomenda-se ler os dois trabalhos em conjunto.

O ensino introdutório de programação envolve dificuldades de compreensão e aplicação dos conceitos, discutidas no mapeamento sistemático de \citeonline{souza2016}. Neste trabalho, o recorte investigado é a exigência de lidar simultaneamente com a construção do algoritmo e com as convenções de escrita da linguagem. Os experimentos de \citeonline{stefik2013} indicam que escolhas de sintaxe afetam o desempenho de iniciantes; esse resultado motiva investigar alternativas de apresentação da linguagem, sem pressupor que uma nova ferramenta produza, por si só, ganhos de aprendizagem.

Em linguagens textuais como C, C++ e Java, o programa precisa respeitar convenções de palavras reservadas, delimitadores e terminadores de instrução para ser aceito pelo tradutor. Um erro nessas convenções pode impedir a tradução antes da execução do algoritmo. A distinção entre a intenção do estudante e a forma aceita pela linguagem é relevante para o ensino, pois parte do esforço inicial recai sobre a identificação de erros sintáticos \cite{aho1995,stefik2013}.

O artefato desenvolvido combina compilação para uma representação intermediária e interpretação dessa representação. Neste texto, \textbf{compilador} designa as etapas que analisam o código-fonte e geram instruções de três endereços; \textbf{interpretador} designa o componente que executa essas instruções. A expressão \textbf{núcleo de tradução e execução} abrange os dois componentes. O sistema é apresentado como um interpretador porque a execução do programa depende desse componente, sem geração de um executável nativo para o usuário. Essa organização é compatível com a distinção entre tradução e execução discutida por \citeonline{aho1995} e \citeonline{nystrom2021}.

O qualificativo \textbf{dinâmico} refere-se à configuração da linguagem recebida a cada nova análise e execução: o usuário pode redefinir palavras-chave, operadores em forma de palavra, literais booleanos e delimitadores, além de selecionar modos de terminação de instruções, de blocos, de tipagem e de arranjos já implementados no núcleo. Essas escolhas não exigem recompilar a aplicação web. O termo não significa que o usuário possa acrescentar produções gramaticais arbitrárias ou alterar a semântica das operações durante a execução de um programa. Os limites dessa configuração são detalhados na Seção~\ref{sec:especificacao-linguagem}.

Diante disso, este trabalho parte da seguinte pergunta: \textbf{em que medida a personalização da sintaxe da linguagem pelo próprio estudante contribui para reduzir a barreira de aprendizagem enfrentada no primeiro contato com a programação?}

Trabalha-se com a hipótese, ainda não submetida a verificação empírica, de que permitir ao aluno configurar os elementos de escrita da linguagem possa favorecer a experimentação de algoritmos e reduzir dificuldades associadas à memorização dessas convenções. A hipótese diz respeito à personalização pelo estudante, inclusive quando os termos escolhidos estão em português, e não apenas à tradução de palavras de uma língua estrangeira.

Cumpre delimitar o alcance deste trabalho diante dessa pergunta. O que aqui se apresenta é a construção do artefato que torna a investigação possível e a delimitação de uma avaliação futura, cujo protocolo ainda precisa ser consolidado. A aplicação desse protocolo junto a estudantes, que permitiria responder empiricamente à pergunta formulada, não integra o escopo deste documento e é registrada como encaminhamento futuro.

\section{Justificativa}

Já existem abordagens que reduzem a exposição inicial a convenções de linguagens profissionais. O Portugol Studio utiliza comandos em português e oferece execução passo a passo e inspeção de variáveis \cite{univali2026}. O Scratch adota programação por blocos e disponibiliza materiais para a criação de projetos interativos \cite{scratch2026}. Portanto, a dificuldade discutida neste trabalho não pode ser atribuída de forma geral ao uso de uma língua estrangeira, nem se presume que as ferramentas existentes ignorem esse problema.

O diferencial investigado é mais delimitado: permitir ao estudante configurar o vocabulário de uma linguagem textual e selecionar variantes de escrita previstas no núcleo, mantendo observáveis os \textit{tokens} e as instruções intermediárias. Isso difere tanto de utilizar um vocabulário em português previamente definido quanto de compor programas por blocos. A comparação restringe-se às ferramentas e aos critérios apresentados na Seção~\ref{sec:trabalhos-relacionados}, sem pretensão de caracterizar a totalidade dos ambientes educacionais existentes.

A proposta é construir esse núcleo configurável e integrá-lo a uma IDE web para experimentação. Uma possível continuidade didática é a resolução de exercícios em linguagens textuais como C++, Java e Python, para as quais o juiz \textit{online} Beecrowd apresenta exemplos de entrada e saída em sua documentação \cite{beecrowd2026faq}. A contribuição da personalização para essa transição é uma expectativa a ser investigada, e não um resultado demonstrado neste documento.

\section{Objetivos}

A seguir são apresentados o objetivo geral e os objetivos específicos que nortearam o desenvolvimento desta pesquisa e a construção da ferramenta.

\subsection{Objetivo Geral}

Projetar e desenvolver um interpretador dinâmico integrado a uma IDE web que permita ao usuário personalizar a sintaxe e o vocabulário da linguagem, com o objetivo de facilitar o processo de ensino-aprendizagem de algoritmos para estudantes iniciantes.

\subsection{Objetivos Específicos}

Para alcançar o objetivo geral, foram definidos os seguintes objetivos específicos:

\begin{itemize}
    \item Especificar a gramática da linguagem e os elementos passíveis de redefinição pelo usuário em tempo de execução;
    \item Projetar os módulos de análise léxica, sintática e semântica com suporte à configuração de vocabulário e à seleção das variantes de linguagem previstas na implementação;
    \item Implementar os módulos projetados, integrando-os em um interpretador dinâmico executável no navegador;
    \item Integrar o núcleo do interpretador à interface web, definindo os contratos de comunicação entre os módulos;
    \item Disponibilizar o núcleo de tradução em ambiente web, verificando seu comportamento por meio de uma suíte de testes automatizados que cubra as construções da linguagem;
    \item Definir o protocolo de avaliação da solução, estabelecendo os critérios que permitirão, em etapa posterior, aferir sua contribuição pedagógica.
\end{itemize}

\section{Estrutura do Trabalho}

Este documento está organizado da seguinte forma: o Capítulo~\ref{cap:referencial} apresenta o referencial teórico e os trabalhos relacionados, incluindo os critérios de seleção das ferramentas. O Capítulo~\ref{cap:proposta} descreve o escopo e o planejamento do trabalho. O Capítulo~\ref{cap:modelagem} detalha a arquitetura e a implementação do núcleo de tradução e execução, a persistência de dados, a integração com a plataforma web e a estratégia de testes. O Capítulo~\ref{cap:consideracoes} retoma os objetivos, as limitações e os encaminhamentos. Os apêndices complementam a descrição com a divisão de responsabilidades, os recursos da API, os testes e o dicionário de dados.
